%  LaTeX support: latex@mdpi.com 
%  For support, please attach all files needed for compiling as well as the log file, and specify your operating system, LaTeX version, and LaTeX editor.

%=================================================================
\documentclass[journal,article,submit,pdftex,moreauthors]{Definitions/mdpi} 
%\documentclass[preprints,article,submit,pdftex,moreauthors]{Definitions/mdpi} 
% For posting an early version of this manuscript as a preprint, you may use "preprints" as the journal. Changing "submit" to "accept" before posting will remove line numbers.

% Below journals will use APA reference format:
% admsci, aieduc, behavsci, businesses, econometrics, economies, education, ejihpe, famsci, games, humans, ijcs, ijfs, journalmedia, jrfm, jsam, languages, peacestud, psycholint, publications, tourismhosp, youth

% Below journals will use Chicago reference format:
% arts, genealogy, histories, humanities, jintelligence, laws, literature, religions, risks, socsci

%--------------------
% Class Options:
%--------------------
%----------
% journal
%----------
% Choose between the following MDPI journals:
% accountaudit, acoustics, actuators, addictions, adhesives, admsci, adolescents, aerobiology, aerospace, agriculture, agriengineering, agrochemicals, agronomy, ai, aichem, aieduc, aieng, aimater, aimed, aipa, air, aisens, algorithms, allergies, alloys, amh, analog, analytica, analytics, anatomia, anesthres, animals, antibiotics, antibodies, antioxidants, applbiosci, appliedchem, appliedmath, appliedphys, applmech, applmicrobiol, applnano, applsci, aquacj, architecture, arm, arthropoda, arts, asc, asi, astronautics, astronomy, atmosphere, atoms, audiolres, automation, axioms, bacteria, batteries, bdcc, behavsci, beverages, biochem, bioengineering, biologics, biology, biomass, biomechanics, biomed, biomedicines, biomedinformatics, biomimetics, biomolecules, biophysica, bioresourbioprod, biosensors, biosphere, biotech, birds, blockchains, bloods, blsf, brainsci, breath, buildings, businesses, cancers, carbon, cardiogenetics, cardiovascmed, catalysts, cells, ceramics, challenges, chemengineering, chemistry, chemosensors, chemproc, children, chips, cimb, civileng, cleantechnol, climate, clinbioenerg, clinpract, clockssleep, cmd, cmtr, coasts, coatings, colloids, colorants, commodities, complexities, complications, compounds, computation, computers, condensedmatter, conservation, constrmater, cosmetics, covid, crops, cryo, cryptography, crystals, csmf, ctn, culture, curroncol, cyber, dairy, data, ddc, dentistry, dermato, dermatopathology, designs, devices, dhi, diabetology, diagnostics, dietetics, digital, disabilities, diseases, diversity, dna, drones, dynamics, earth, ebj, ecm, ecologies, econometrics, economies, edm, education, eesp, ejihpe, electricity, electrochem, electronicmat, electronics, encyclopedia, endocrines, energies, eng, engproc, entomology, entropy, environments, environremediat, epidemiologia, epigenomes, esa, est, famsci, fermentation, fibers, fintech, fire, fishes, fluids, foods, forecasting, forensicsci, forests, fossstud, foundations, fractalfract, fuels, future, futureinternet, futurepharmacol, futurephys, futuretransp, galaxies, games, gases, gastroent, gastrointestdisord, gastronomy, gels, genealogy, genes, geographies, geohazards, geomatics, geometry, geosciences, geotechnics, geriatrics, germs, glacies, grasses, green, greenhealth, gucdd, hardware, hazardousmatters, healthcare, hearts, hemato, hematolrep, hep, heritage, higheredu, highthroughput, histories, horticulturae, hospitals, humanities, humans, hydrobiology, hydrogen, hydrology, hydropower, hygiene, idr, iic, ijcs, ijem, ijerph, ijfs, ijgi, ijmd, ijms, ijns, ijom, ijpb, ijt, ijtm, ijtpp, ime, immuno, informatics, information, infrastructures, inorganics, insects, instruments, inventions, iot, j, jaestheticmed, jal, jcdd, jcm, jcp, jcrm, jcs, jcto, jdad, jdb, jdream, jemr, jeta, jfb, jfmk, jgbg, jgg, jimaging, jintelligence, jlpea, jmahp, jmmp, jmms, jmp, jmse, jne, jnt, jof, joi, joitmc, joma, jop, joptm, jor, journalmedia, jox, jpbi, jphytomed, jpm, jrfm, jsam, jsan, jtaer, jvd, jzbg, kidneydial, kinasesphosphatases, knowledge, labmed, laboratories, lae, land, languages, laws, life, lights, limnolrev, lipidology, liquids, literature, livers, logics, logistics, lubricants, lymphatics, machines, macromol, magnetism, magnetochemistry, make, marinedrugs, materials, materproc, mathematics, mca, measurements, medicina, medicines, medsci, membranes, merits, metabolites, metals, meteorology, methane, metrics, metrology, micro, microarrays, microbiolres, microelectronics, micromachines, microorganisms, microplastics, microwave, minerals, mining, mmphys, modelling, molbank, molecules, mps, msf, mti, multimedia, muscles, nanoenergyadv, nanomanufacturing, nanomaterials, ncrna, ndt, network, neuroglia, neuroimaging, neurolint, neurosci, nitrogen, notspecified, nursrep, nutraceuticals, nutrients, obesities, occuphealth, oceans, ohbm, onco, optics, oral, organics, organoids, osteology, oxygen, pandemics, parasites, parasitologia, particles, pathogens, pathophysiology, peacestud, pediatrrep, pets, pharmaceuticals, pharmaceutics, pharmacoepidemiology, pharmacy, philosophies, photochem, photonics, phycology, physchem, physics, physiologia, plants, plasma, platforms, pollutants, polymers, polysaccharides, populations, poultry, powders, precipitation, precisoncol, preprints, proceedings, processes, prosthesis, proteomes, psf, psychiatryint, psychoactives, psycholint, publications, purification, quantumrep, quaternary, qubs, radiation, rdt, reactions, realestate, receptors, recycling, regeneration, religions, remotesensing, reports, reprodmed, resources, rheumato, risks, rjpm, robotics, rsee, ruminants, safety, sci, scipharm, sclerosis, seeds, sensors, separations, sexes, shi, signals, sinusitis, siuj, skins, smartcities, sna, societies, socsci, software, soilsystems, solar, solids, spectroscj, sports, standards, stats, std, stratsediment, stresses, surfaces, surgeries, suschem, sustainability, symmetry, synbio, systems, tae, targets, taxonomy, technologies, telecom, test, textiles, thalassrep, therapeutics, thermo, timespace, tomography, tourismhosp, toxics, toxins, tph, transplantology, transportation, traumacare, traumas, tri, tropicalmed, universe, urbansci, uro, vaccines, vehicles, venereology, vetsci, vibration, virtualworlds, viruses, vision, waste, water, welding, wem, wevj, wild, wind, women, world, youth, zoonoticdis

%---------
% article
%---------
% The default type of manuscript is "article", but can be replaced by: 
% abstract, addendum, article, benchmark, book, bookreview, briefcommunication, briefreport, casereport, changes, clinicopathologicalchallenge, comment, commentary, communication, conceptpaper, conferenceproceedings, correction, conferencereport, creative, datadescriptor, discussion, entry, expressionofconcern, extendedabstract, editorial, essay, erratum, fieldguide, hypothesis, interestingimages, letter, meetingreport, monograph, newbookreceived, obituary, opinion, proceedingpaper, projectreport, reply, retraction, review, perspective, protocol, shortnote, studyprotocol, supfile, systematicreview, technicalnote, viewpoint, guidelines, registeredreport, tutorial,  giantsinurology, urologyaroundtheworld
% supfile = supplementary materials

%----------
% submit
%----------
% The class option "submit" will be changed to "accept" by the Editorial Office when the paper is accepted. This will only make changes to the frontpage (e.g., the logo of the journal will get visible), the headings, and the copyright information. Also, line numbering will be removed. Journal info and pagination for accepted papers will also be assigned by the Editorial Office.

%------------------
% moreauthors
%------------------
% If there is only one author the class option oneauthor should be used. Otherwise use the class option moreauthors.

%---------
% pdftex
%---------
% The option pdftex is for use with pdfLaTeX. Remove "pdftex" for (1) compiling with LaTeX & dvi2pdf (if eps figures are used) or for (2) compiling with XeLaTeX.

%=================================================================
% MDPI internal commands - do not modify
\firstpage{1} 
\makeatletter 
\setcounter{page}{\@firstpage} 
\makeatother
\pubvolume{1}
\issuenum{1}
\articlenumber{0}
\pubyear{2026}
\copyrightyear{2025}
%\externaleditor{Firstname Lastname} % More than 1 editor, please add `` and '' before the last editor name
\datereceived{ } 
\daterevised{ } % Comment out if no revised date
\dateaccepted{ } 
\datepublished{ } 
%\datecorrected{} % For corrected papers: "Corrected: XXX" date in the original paper.
%\dateretracted{} % For retracted papers: "Retracted: XXX" date in the original paper.
\hreflink{https://doi.org/} % If needed use \linebreak
%\doinum{}
%\pdfoutput=1 % Uncommented for upload to arXiv.org
%\CorrStatement{yes}  % For updates
%\longauthorlist{yes} % For many authors that exceed the left citation part
%\IsAssociation{yes} % For association journals

%=================================================================
% Add packages and commands here. The following packages are loaded in our class file: fontenc, inputenc, calc, indentfirst, fancyhdr, graphicx, epstopdf, lastpage, ifthen, float, amsmath, amssymb, lineno, setspace, enumitem, mathpazo, booktabs, titlesec, etoolbox, tabto, xcolor, colortbl, soul, multirow, microtype, tikz, totcount, changepage, attrib, upgreek, array, tabularx, pbox, ragged2e, tocloft, marginnote, marginfix, enotez, amsthm, natbib, hyperref, cleveref, scrextend, url, geometry, newfloat, caption, draftwatermark, seqsplit
% cleveref: load \crefname definitions after \begin{document}

%=================================================================
% Please use the following mathematics environments: Theorem, Lemma, Corollary, Proposition, Characterization, Property, Problem, Example, ExamplesandDefinitions, Hypothesis, Remark, Definition, Notation, Assumption
%% For proofs, please use the proof environment (the amsthm package is loaded by the MDPI class).

%=================================================================
% Full title of the paper (Capitalized)
\Title{LOW COST GAN FOR DEPTH ESTIMATION}

% Author Orchid ID: enter ID or remove command
\newcommand{\orcidauthorA}{0000-0000-0000-000X} % Add \orcidA{} behind the author's name
%\newcommand{\orcidauthorB}{0000-0000-0000-000X} % Add \orcidB{} behind the author's name

% Authors, for the paper (add full first names)
\Author{Arif Suganda $^{1}$\orcidA{}}

%\longauthorlist{yes}

% MDPI internal command: Authors, for metadata in PDF
\AuthorNames{Arif Suganda}

% Affiliations / Addresses (Add [1] after \address if there is only one affiliation.)
\address{%
$^{1}$ \quad Department of Technology Information, Tadulako University, Palu, Indonesia; f5522530002@untad.ac.id}

% Contact information of the corresponding author
\corres{Correspondence: arifsuganda.su@gmail.com}

% Abstract (Do not insert blank lines, i.e. \\) 
\abstract{Depth estimation, one of the leading research in computer vision, from autonomous driving to 3D modeling, has play a key role to artificial perception of the machine. This paper present reconstruction of indoor 3D plane focusing on optimal cost to accuracy ratio, using simple Generative Adversarial Networks (GANs) trained using NYU Depth v2 dataset for indoor plane. Using this limitation our approach in lowering cost for optimal cost to accuracy ratio has achieved xx\% accuracy using only xx training time resulted in xx times yield compared to other approach.}

% Keywords
\keyword{monocular depth estimation; GAN; indoor environment} 

\begin{document}

\section{Introduction}
Monocular depth estimation (MDE) is a classic ill\mbox{-}posed problem for using only a single RGB image to predict full image depth without spatial information.
Despite it's challenge, MDE is still the most popular research topic for depth estimation~\cite{guizilini9156708} for many different poblem, 
ranging from robotic~\cite{chang9010976}, autonomous driving, to 3d modeling.
There are many approach explored by the current researches in MDE, from optimizing hardware~\cite{chang9010976}, improving dataset~\cite{guizilini9156708}, and/or using new set of architecture~\cite{guizilini9156708,hendra10113647}.
Through these approaches, they share the same commonality, that is the current research landscape on MDE heavily dominated by deep leading (DL) models~\cite{masomians22145353}.

Since Eigen~\cite{eigen14062283}, MDE has been replicated into many different research approaches with varied results and still posed the same challenge of capturing global dan local depth information from a single image.
These varied approaches had consistently achieve 90\% or above accuracy~\cite{khan20082272} using many different architecture model, dataset, and even modified camera to help capture image depth.
The high accuracy result comes at the high cost of creating data with ground truth and high cost of supervised training~\cite{zhao101007,masomians22145353,khan20082272},
Which limit their practicality and applicabilty in real-time and resource-constrained environment~\cite{zhao101007,khan20082272}

To reduce this reliance on ground truth data and supervised training method, a self-supervised and semi-supervised training method have been proposed,
leveraging on photometric and geometric consistency from stereo image pairs or monocular video sequences~\cite{Zhou2017Unsupervised,Godard2017Monodepth}.
Although these approaches reduce data labeling, their performace often degrade in textureless area and complex indoor environment, which is common in real-world application~\cite{masomians22145353}.
Moreover, most of the existing method formulate image depth as a pixel-wise regression task, which frequently lead to over-smoothed depth predictions and insufficient preservation of global scene structure, 
particularly planar surfaces and sharp depth discontinuities~\cite{khan20082272}.

Generative Adversarial Networks (GANs) recently has also enter as one of the alternative to mitigated this limitations. 
By incorporating adversarial learning, GAN-based approaches use the model to align predicted depth maps with the distribution of real depth data, thereby, improving gobal consistency and structural realism~\cite{Chen2019GANDepth,Isola2017Pix2Pix,Khan2023GANDepth}.
Several studies have proved improvement in edge preservation and percetual quality compared to the conventional CNN-based methods
However, these gains are commonly achieved using a heavy backbone networks which increate the training and inference cost.
As highlighted in recent reviews, most GAN-based approaches prioritize depth accuracy or visual quality, largely overlooking computational efficiency [2].

This research proposed a cost-effective GAN-based MDE framework, designed for indoor environment to be posed as a constraint for 3D plan reconstruction.
The proposed approach employs lightweight generator networks optimized for efficiency, while the adversarial discriminator enforces global depth structure and plane consistency.
The method is evaluated in the NYU Depth V2 dataset~\cite{Silberman2012NYU}, a widely used benchmark for indoor depth estimation.
The estimated depth maps are further assessed through a 3D plane reconstruction, to demonstrate the practical applicability of the proposed approach in indoor scene understanding.

The main contribution of this research can be summarized as follows:
\begin{enumerate}
    \item A lightweight GAN-based MDE framework designed to achieve a favorable trade-off between accuracy and computational cost.
    \item An adversarial learning strategy to improve global structure consistency and planar depth quality without relying on heavy networks backbone.
    \item A Comprehensive evaluation of indoor 3D plane reconstruction performance on NYU Depth V2.
\end{enumerate}

The remainder of this paper is organized as follows. Section 2 reviews related work on deep learning–based and GAN-based monocular depth estimation methods. 
Section 3 describes the proposed methodology in detail. Section 4 presents the experimental results and it's comparative evaluations. 
Finally, Section 5 concludes the paper and discusses future research directions.
%%%%%%%%%%%%%%%%%%%%%%%%%%%%%%%%%%%%%%%%%%
\section{Related Work}

\subsection{CNN-Based Monocular Depth Estimation}
Monocular depth estimation has been implemented with various approaches, starting from Eigen~\cite{eigen14062283} with the first CNN implementation of MDE 
which introduced a multi-scale architecture to capture global scene layout and local depth details. 
Subsequent work complement this framework using geometry cues such as nomal and semantic labels using shared multi-scale representation~\cite{lee2021bigsmallmultiscalelocal}.

Later studies fully convolutional encoder-decoder architectures to improve spatial resolution and convergence behavior using with skipped connections and residual learning~\cite{zhang2019patternaffinitivepropagationdepthsurface}.
These method achieved a strong performance on indoor benchmark such as NYU Depth V2~\cite{zhou2017unsupervisedlearningdepthegomotion}. 
However, these models typically rely on deep backbone network with high computational and memory requirements, which limits their suitability for efficiency-critical applications.

\subsection{Efficiency-Oriented Depth Estimation Methods}
Several studies have explored a lightweight MDE model, such as using an efficient and lightweight encoder-decoder network architecture~\cite{Wofk2019FastDepth},
use combination of shallow encoder–decoder CNNs and hybrid loss at the cost of degraded performance on a high complex problem~\cite{dong2021mobilexnetefficientconvolutionalneural}.
and other approaches using simplified networks or distillation strategy to prune the networks reducing the inference cost.

Despite improved efficiency, lightweight models often rely on pixel-wise loss functions and reduced network capability. 
This process often leads to over-smoothed depth maps and degraded structural accuracy, which is particularly degrading in indoor environments dominated by planar surfaces~\cite{khan20082272,masomians22145353}.

\subsection{GAN-Based Monocular Depth Estimation}
Generative adversarial networks (GANs) have been introduced to MDE to address the limitations of regression-based models, particularly over-smoothed depth maps and degraded structural accuracy.
Conditional GAN (cGans) learn mapping between RGB images to depth maps while enforcing the distribution-level similarity through supervision of the adversarial model~\cite{Hao2023GAN}.

\section{Proposed Method}
In this section we introduce our proposed method for MDE using a simple GANs architecture,
first we explain our network architecture, then the evaluation method used to our method.

\subsection{Network Architecture}
Our proposed method employs simple GANs model using a generator and a indiscriminator model demonstrated in figure~\ref{fig:gan}.
The proposed method employs a generator networks to predict the depth maps from RGB images, while adversarial discriminator enforces global structure and plane consistency in the predicted depth.
\begin{figure*}[t]\label{fig:gan}
    \centering
    \includegraphics[width=0.8\textwidth]{gan.png}
    \caption{GANs architecture}
\end{figure*}

\subsection{Evaluation}
Our method will be evaluated using the same benchmark for NYU Depth V2 datasets, using these performace metrics~\ref{eq:rmse},~\ref{eq:rel}, and~\ref{eq:accuracy_ut}:
\begin{enumerate}
    \item Root Mean Squared Error (RMSE):
    \begin{equation}\label{eq:rmse}
        \mathrm{RMSE} = \sqrt{\frac{1}{N}\sum_{i=1}^{N} \left( y_i - \hat{y}_i \right)^2}
    \end{equation}
    \item Average Relative Error (REL):
    \begin{equation}\label{eq:rel}
        \mathrm{REL} = \frac{1}{N}\sum_{i=1}^{N} \frac{\left| y_i - \hat{y}_i \right|}{y}
    \end{equation}
    \item Accuracy under Threshold ($\delta i$):
    \begin{equation}\label{eq:accuracy_ut}
        \delta_T = \frac{1}{N}
        \left|
        \left\{
        i \;\middle|\;
        \max\left(
        \frac{\hat{y}_i}{y_i},
        \frac{y_i}{\hat{y}_i}
        \right) < T
        \right\}
        \right|
    \end{equation}
\end{enumerate}


\section{Experiment}

\section{Discussion}

\section{Conclusion}

\reftitle{References}

\bibliography{refs}

\PublishersNote{}

\end{document}

